\section{External Consistency Check: CADE Cartel Convictions}
\label{sec:cade}

Before examining prices, we check whether FL firms appear in
environments linked to known cartels. CADE (Conselho Administrativo
de Defesa Econ\^omica) convicted 65 procurement cartel cases during
2009--2019; 47 convicted firms match to BEC via CNPJ. Among 2,735 FL
firms, 193 (7.1\%) co-participate with at least one CADE-convicted
firm---3.5 times the expected rate under a permutation null stratified
by participation quartile ($p < 0.001$; 1,000 iterations). The 92.9\%
without CADE overlap reflects the well-documented low base rate of
cartel prosecution \citep{connor2007price}: CADE convictions are the
visible tip; the screen is designed to flag environments that may include
undetected collusion.

\paragraph{Interpretation.}
Co-participation measures market proximity, not cartel membership.
Once we condition on participation count, the FL--CADE overlap rate
becomes indistinguishable from that of other always-losers ($p = 0.93$;
Section~\ref{sec:ground_truth}): the excess is driven by high
participation volume. This is the screen's operating principle---a
participation-based triage tool flags firms that participate heavily in
suspicious markets, without requiring the researcher to distinguish
proximity from direct involvement. The permutation test
(3.5$\times$, stratified by participation quartile) confirms that the
excess is real, not mechanical.

\paragraph{Dual pattern.}
Most CADE-convicted firms in BEC are \emph{winners}---the designated
recipients of cartel rents. But three convicted firms are themselves
classified as frequent losers (Table~\ref{tab:cade_fl_firms}). All
three participated across multiple item groups and years without ever
winning---the defining FL pattern.

\begin{table}[htbp]
\centering
\caption{CADE-convicted firms classified as frequent losers}
\label{tab:cade_fl_firms}
\small
\begin{tabular}{lllcc}
\toprule
Firm & Cartel & CADE process & Tenders & Wins \\
\midrule
Sol Tecnologia & Solar water heaters & 08012.001273/2010-24 & 84 & 0 \\
Nova Esperan\c{c}a Locadora & School transportation & 08700.005876/2019-85 & 97 & 0 \\
Jofran Com\'ercio & Trash bags (Op.\ Colludium) & 08700.005789/2015-02 & 65 & 0 \\
\bottomrule
\end{tabular}
\end{table}

\paragraph{Co-bidder analysis.}
FL firms appear in 5.6\% of CADE-firm tenders versus a 2.1\% baseline
(2.6$\times$). For individual convicted firms the rate is much
higher---up to 69\% for Mayfran (school transportation). Excluding
the three directly convicted FL firms and their 246 tender-items
leaves $\hat{\beta}$ unchanged at 0.064.

\paragraph{Beyond known cartels.}
Dropping all 31,447 CADE-involved tender-items yields
$\hat{\beta} = 0.062$ ($N = 1{,}622{,}954$), barely changed from the
full-sample 0.064 (Table~\ref{tab:excl_cade}). The FL--price
association is not driven by tenders CADE has already flagged---the
screen reaches markets that fall outside the enforcement record.

